Put \(p=\pi_1\) and \(q=\pi'_1\). By \texttt{ass:source-overlap}, \(0<p<1\); hence every denominator below is nonzero. By \texttt{ass:source-product-design}, the two assignments have source probabilities
\[
p_{\pi}(0)=1-p,
\qquad
p_{\pi}(1)=p.
\]
By \texttt{ass:target-product-policy} and \texttt{def:local-target-mass}, their target local-cell probabilities are \(p_{\pi',1}(0)=1-q\) and \(p_{\pi',1}(1)=q\).

We first identify the parameter set in this special case. The hypotheses give \(X_1=\{0,1\}\) and \(\mathcal S_1=\{\varnothing,\{1\}\}\). By \texttt{def:local-fourier-character},
\[
\chi_{1,\varnothing}^{\pi}(x)=1,
\qquad
\chi_{1,\{1\}}^{\pi}(x)=\frac{x-p}{\sqrt{p(1-p)}}.
\]
The second character takes different values at \(0\) and \(1\), since their difference is \(1/\sqrt{p(1-p)}\ne0\). Thus these two characters are linearly independent. Because \(\mathbb R^{X_1}\) has dimension two, \texttt{def:local-fourier-space} yields \(H_1=\mathbb R^{X_1}\). It then follows from \texttt{def:bounded-fourier-class} that
\[
K=[-1,1]^{X_1}=[-1,1]^2.
\]
Write a schedule as \(v=(v_0,v_1)\), where \(v_x=v_1(x)\). By \texttt{def:observation-operator}, assignment \(z=x\) reveals \(N_xv=(v_x)\). By \texttt{def:target-value} and the target-cell probabilities above,
\[
\theta(v)=(1-q)v_0+qv_1.
\]

For the upper bound, choose the assignment-indexed linear weights \(a_0=1-q\) and \(a_1=q\). The resulting rule is
\[
\delta_a(0,N_0v)=(1-q)v_0,
\qquad
\delta_a(1,N_1v)=qv_1.
\]
Consequently its two assignment-specific errors are
\[
\delta_a(0,N_0v)-\theta(v)=-qv_1,
\qquad
\delta_a(1,N_1v)-\theta(v)=-(1-q)v_0.
\]
Its design risk therefore satisfies
\[
\begin{aligned}
r_a(v)
&=(1-p)q^2v_1^2+p(1-q)^2v_0^2\\
&\le (1-p)q^2+p(1-q)^2=:A,
\end{aligned}
\]
where the inequality uses \(|v_0|,|v_1|\le1\). Equality holds, for example, at \((v_0,v_1)=(1,1)\). Hence \texttt{def:linear-frontier} gives \(F\le A\), and this displayed rule attains worst-case risk \(A\).

It remains to prove that no measurable nonlinear rule does better. Fix an arbitrary Borel rule \(\delta\), and abbreviate \(d_x(y)=\delta(x,(y))\). Average its risk over the genuine probability distribution under which \(V_0\) and \(V_1\) are independent and each is uniform on \(\{-1,1\}\). This distribution is supported on \(K\). Conditional on assignment zero, set \(T_0=d_0(V_0)-(1-q)V_0\). Independence, \(\mathbb E[V_1]=0\), and \(V_1^2=1\) give
\[
\begin{aligned}
\mathbb E\!\left[\bigl(d_0(V_0)-\theta(V)\bigr)^2\right]
&=\mathbb E\!\left[(T_0-qV_1)^2\right]\\
&=\mathbb E[T_0^2]-2q\,\mathbb E[T_0]\mathbb E[V_1]+q^2\mathbb E[V_1^2]\\
&=\mathbb E[T_0^2]+q^2\\
&\ge q^2.
\end{aligned}
\]
Similarly, conditional on assignment one, put \(T_1=d_1(V_1)-qV_1\). Independence, \(\mathbb E[V_0]=0\), and \(V_0^2=1\) imply
\[
\mathbb E\!\left[\bigl(d_1(V_1)-\theta(V)\bigr)^2\right]
=\mathbb E[T_1^2]+(1-q)^2
\ge(1-q)^2.
\]
Combining the two assignments with their source probabilities shows that the average, over this four-point prior, of the design risk of \(\delta\) is at least
\[
(1-p)q^2+p(1-q)^2=A.
\]
The supremum of a function over \(K\) is at least its average under any probability distribution supported on \(K\). Therefore every Borel \(\delta\) has worst-case risk at least \(A\), and \texttt{def:unrestricted-frontier} yields \(R^{\star}\ge A\). Conversely, every assignment-indexed linear rule is Borel, so the two infima in \texttt{def:linear-frontier} and \texttt{def:unrestricted-frontier} satisfy \(R^{\star}\le F\). Together with \(F\le A\),
\[
A\le R^{\star}\le F\le A,
\]
which proves
\[
R^{\star}=F=(1-\pi_1)(\pi'_1)^2+\pi_1(1-\pi'_1)^2.
\]
This argument explicitly enumerates both assignments and all four support points of the lower-bound prior, so it also supplies the claimed smallest-instance stress test.

Finally, suppose \(v_1(0)=v_1(1)=c\). Boundedness gives \(|c|\le1\), and the target formula gives \(\theta(v)=(1-q)c+qc=c\). Under either \(z=0\) or \(z=1\), \texttt{def:observation-operator} gives \((N_zv)_1=c\). Thus the Borel rule \(\delta(z,N_zv)=(N_zv)_1\) equals \(\theta(v)\) under each assignment, and its squared error is identically zero.